% !TEX program = pdflatex
\documentclass[11pt]{article}
\usepackage[margin=1in]{geometry}
\usepackage{amsmath,amssymb,amsthm}
\usepackage{graphicx}
\usepackage{booktabs}
\usepackage{siunitx}
\usepackage{hyperref}
\usepackage{natbib}
\usepackage{enumitem}
\usepackage{listings}
\usepackage{float}
\usepackage{xcolor}
\lstset{basicstyle=\ttfamily\small,breaklines=true,frame=single}
\setlist{noitemsep,topsep=2pt}

\title{Combining BMI and Electrocardiogram Features for Hypertension Risk Prediction: A Pragmatic Development and Validation Protocol}
\author{Ryan O. van Gelder \\ \\ Citizen Gardens}
\date{October 23, 2025}

\begin{document}
\maketitle

\begin{abstract}
We present a production-oriented protocol to develop and validate a model that estimates baseline or one-year incident hypertension using body mass index (BMI) and electrocardiogram (EKG) features. The design emphasizes clear baselines, rigorous calibration, clinical utility analysis, and deployment safeguards. We include reference implementations for feature extraction and modeling.\footnote{All code snippets are illustrative and assume standard Python packages.}
\end{abstract}

\section{Objective}
Let $y\in\{0,1\}$ denote prevalent hypertension at index or incident hypertension within 1 year. Given features $x$ from BMI, demographics, and EKG, estimate $\hat p=\Pr(y=1\mid x)$.

\section{Cohort and Labels}
\begin{itemize}
  \item Adults with paired BMI and EKG within $\leq$ 30 days.
  \item Prevalent HTN: diagnosis codes, antihypertensive meds, or blood pressure $\geq$140/90 on two dates.
  \item Incident HTN: no HTN at index; HTN within $365\pm30$ days.
  \item Exclusions: AFib during recording, paced rhythm, missing BMI, unreadable EKG.
  \item Sensitivity: evaluate 7/14/30-day pairing windows and report cohort yield and metric stability.
  \item Report metrics separately for prevalent and incident HTN.
\end{itemize}

\section{Data Schema and Notation}
Demographics (age, sex); BMI (kg/m$^2$); EKG (12- or single-lead, $f_s\geq 250$ Hz, $\geq$10 s); metadata (device, site, date). Raw signal per lead $s_\ell(t)$ for $\ell=1..L$. R-peak times $\{r_i\}$; NN intervals $NN_i=r_{i+1}-r_i$; heart rate series $HR_i=60/NN_i$.

\section{Signal Processing}
Bandpass 0.5--40 Hz; notch if needed. Detect R-peaks with Pan--Tompkins or comparable \citep{pan1985real}. Manual spot check on 1\%. Artifact rules: drop ectopic beats; require $\geq$30 NN intervals. Targeted manual review on 5\% of low-SNR or irregular cases, escalating to 10\% if R-peak error $>$1\%. Quality gates: SNR $\geq$6 dB and $<20\%$ beats rejected.

\section{Features}
\subsection{Time-domain HRV}
Mean NN, SDNN, RMSSD, pNN50 per Task Force standards \citep{taskforce1996hrv}. For example, $\mathrm{SDNN}=\sqrt{\tfrac{1}{n-1}\sum (NN_i-\overline{NN})^2}$.
\subsection{Frequency-domain HRV}
Welch PSD \citep{welch1967}, bands: VLF 0.003--0.04 Hz, LF 0.04--0.15 Hz, HF 0.15--0.4 Hz; features: LF, HF, LF/HF, total power \citep{taskforce1996hrv}.
\subsection{Morphology}
Intervals PR, QRS, QT and QTc using Bazett and Fridericia \citep{bazett1920interval,fridericia1920systole}. Amplitudes (R,S,T), ST deviation at J+60 ms, T-wave axis; optional PCA for vector features.
\subsection{Quality}
SNR estimate, percent beats rejected, missing-lead flags.

\section{Models}
\subsection{Baseline A (Clinical)}
Logistic regression: $\mathrm{logit}(\hat p)=\beta_0+\beta_1\,\mathrm{BMI}+\beta_2\,\mathrm{Age}+\beta_3\,\mathrm{Sex}$.
\subsection{Baseline B (EKG-only)}
Regularized logistic regression or gradient boosting on EKG features.
\subsection{Model C (Combined)}
Concatenate Baseline A covariates with EKG features; same learner as B.
\subsection{Optional Deep Model}
A 1D-CNN on raw lead II: three Conv--BN--ReLU blocks (kernel 7, stride 2), maxpool, global average pooling, dense head; concatenate BMI, age, sex at the penultimate layer. Binary cross-entropy with class weighting. Research-only unless it outperforms Model C and meets calibration and explainability parity.

\section{Training}
Stratified 70/15/15 split by patient; hold out an external site if available. Standardize continuous features on train folds. Median impute with missing-indicators. Class-imbalance handled by inverse prevalence weights or focal loss for deep models. Hyperparameters via 5-fold CV on training data only.

\section{Evaluation}
Discrimination: AUROC and AUPRC. Calibration: Brier score \citep{brier1950}, expected calibration error, reliability plots \citep{steyerberg2019clinical}. Clinical utility: decision curve analysis \citep{vickers2006dca}; number needed to evaluate (NNE). Compare AUROC with DeLong's test \citep{delong1988}. Report $\Delta$ vs baselines with CIs and $p$-values. Subgroups: prevalent vs incident; by device and site. Net reclassification improvement vs Baseline A at 5\%, 10\%, 20\% thresholds \citep{pencina2008nri}. Decision curve net benefit with 95\% CIs over 1--30\%.

\section{Explainability}
Linear/GBM: standardized coefficients or feature gains; SHAP summaries and dependence plots for top features \citep{lundberg2017shap}. CNN: Grad-CAM on the last conv layer \citep{selvaraju2017gradcam}; require physiological plausibility.

\section{Ablations}
Remove BMI from Model C; remove EKG features from Model C; remove HRV or morphology blocks; evaluate 10 s vs 60 s EKG if available.

\section{External Validation}
Train on Site A, test on Site B. Report domain shift (prevalence, devices, demographics). If drift is detected, recalibrate with Platt scaling \citep{platt1999} or isotonic regression \citep{zadrozny2002calibration} using 10\% of Site B; report pre/post metrics.

\section{Thresholding}
Select operating points by Youden's $J$ \citep{youden1950index} and by fixed sensitivity 0.80. Report PPV, NPV, LR$+$, LR$-$, and confusion matrices per site. Governance: register owner and review cadence; revisit threshold when prevalence shifts by $>$25\% relative or ECE $>$0.05.

\section{Sample Size Guidance}
For logistic regression with $k$ predictors, target $\ge 20k$ events \citep{steyerberg2019clinical}. For AUROC CI width $\le 0.05$ at prevalence $\approx0.3$, aim for $n\approx$2{,}000--5{,}000 with $\ge$600 events. Conduct a feasibility check under 7/14/30-day pairing windows; target $\ge$20k events per 1{,}000 candidate predictors after exclusions \citep{riley2020sample}.

\section{Reproducibility}
Deterministic pipeline with seeds. Version data slices, code, and hyperparameters. Provide a model card with data sources, inclusion/exclusion, metrics with CIs, calibration, failure modes, and fairness slices (TRIPOD-aligned reporting \citep{collins2015tripod}).

\section{Deployment and Monitoring}
Export model as ONNX or PMML and include the preprocessing graph. Runtime checks for input ranges, lead presence, and signal quality. Online monitoring: prevalence, AUROC proxy via risk ranking, calibration via periodic outcome linkage, population stability index (PSI), and net benefit at the chosen threshold. Recalibrate quarterly or when ECE $>$0.05.

\section{Reference Implementations}
\subsection{HRV and morphology feature extraction (Python)}
\begin{lstlisting}[language=Python]
import numpy as np
from scipy.signal import butter, filtfilt, welch

FS = 500  # sampling rate

def bandpass(signal, fs=FS, lo=0.5, hi=40.0, order=4):
    b, a = butter(order, [lo/(fs/2), hi/(fs/2)], btype='band')
    return filtfilt(b, a, signal)

# r_peaks assumed provided by a Pan–Tompkins-like detector

def nn_intervals(r_peaks, fs=FS):
    rr = np.diff(r_peaks) / fs
    mask = (rr > 0.3) & (rr < 2.0)  # drop ectopy/outliers
    nn = rr[mask]
    return nn

def time_domain_hrv(nn):
    mean_nn = np.mean(nn)
    sdnn = np.std(nn, ddof=1)
    rmssd = np.sqrt(np.mean(np.diff(nn)**2))
    diff50 = np.sum(np.abs(np.diff(nn)) > 0.05)
    pnn50 = diff50 / max(len(nn)-1, 1)
    return {
        'mean_nn': mean_nn,
        'sdnn': sdnn,
        'rmssd': rmssd,
        'pnn50': pnn50,
    }

def frequency_domain_hrv(nn, fs_nn=4.0):
    # interpolate NN series at 4 Hz, then Welch PSD
    t = np.cumsum(nn)
    t_uniform = np.arange(0, t[-1], 1/fs_nn)
    nn_interp = np.interp(t_uniform, t[:-1], nn[:-1])
    f, pxx = welch(nn_interp, fs=fs_nn, nperseg=256)
    def band_power(lo, hi):
        idx = (f >= lo) & (f < hi)
        return np.trapz(pxx[idx], f[idx])
    vlf = band_power(0.003, 0.04)
    lf  = band_power(0.04, 0.15)
    hf  = band_power(0.15, 0.40)
    total = vlf + lf + hf
    return {'vlf': vlf, 'lf': lf, 'hf': hf, 'lf_hf': lf/(hf+1e-9), 'total': total}
\end{lstlisting}

\subsection{Classical pipeline (scikit-learn)}
\begin{lstlisting}[language=Python]
import numpy as np
import pandas as pd
from sklearn.compose import ColumnTransformer
from sklearn.preprocessing import StandardScaler
from sklearn.impute import SimpleImputer
from sklearn.linear_model import LogisticRegression
from sklearn.metrics import roc_auc_score, average_precision_score
from sklearn.pipeline import Pipeline

num_cols = ['BMI', 'Age'] + ekg_numeric_features
cat_cols = ['Sex']

preprocess = ColumnTransformer([
    ('num', Pipeline([
        ('imp', SimpleImputer(strategy='median')),
        ('sc', StandardScaler())
    ]), num_cols),
    ('cat', SimpleImputer(strategy='most_frequent'), cat_cols)
])

clf = Pipeline([
    ('prep', preprocess),
    ('logit', LogisticRegression(max_iter=2000, class_weight='balanced', solver='liblinear'))
])

clf.fit(X_train, y_train)
proba = clf.predict_proba(X_valid)[:,1]
print('AUROC', roc_auc_score(y_valid, proba))
print('AUPRC', average_precision_score(y_valid, proba))
\end{lstlisting}

\subsection{1D-CNN for raw EKG (Keras)}
\begin{lstlisting}[language=Python]
import tensorflow as tf
from tensorflow.keras import layers, models

# x_sig: (N, T, 1) raw lead II; x_tab: (N, 3) BMI, age, sex
inp_sig = layers.Input(shape=(T, 1))
x = layers.Conv1D(32, 7, strides=2, padding='same', activation='relu')(inp_sig)
x = layers.BatchNormalization()(x)
x = layers.Conv1D(64, 7, strides=2, padding='same', activation='relu')(x)
x = layers.BatchNormalization()(x)
x = layers.Conv1D(128,7, strides=2, padding='same', activation='relu')(x)
x = layers.BatchNormalization()(x)
x = layers.MaxPooling1D()(x)
x = layers.GlobalAveragePooling1D()(x)

inp_tab = layers.Input(shape=(3,))
conc = layers.Concatenate()([x, inp_tab])
d = layers.Dense(64, activation='relu')(conc)
out = layers.Dense(1, activation='sigmoid')(d)

model = models.Model([inp_sig, inp_tab], out)
model.compile(optimizer=tf.keras.optimizers.Adam(1e-3),
              loss='binary_crossentropy',
              metrics=[tf.keras.metrics.AUC(name='auroc'),
                       tf.keras.metrics.AUC(name='auprc', curve='PR')])
# model.fit([Xsig_train, Xtab_train], y_train, validation_data=([Xsig_val, Xtab_val], y_val),
#           epochs=20, batch_size=64, class_weight={0:1,1:w})
\end{lstlisting}

\section{Limitations}
Cohort selection may reduce generalizability; signal quality and device heterogeneity can degrade performance; interpretability constraints for deep learning; prospective and interventional validation are out of scope.

\section*{Acknowledgments}
None.

\bibliographystyle{unsrtnat}
\begin{thebibliography}{99}
\bibitem[Pan and Tompkins(1985)]{pan1985real} J.~Pan and W.~J. Tompkins. A real-time QRS detection algorithm. \emph{IEEE Trans. Biomed. Eng.}, 32(3):230--236, 1985.
\bibitem[Task Force(1996)]{taskforce1996hrv} Task Force of the ESC and NASPE. Heart rate variability: standards of measurement, physiological interpretation, and clinical use. \emph{Circulation}, 93(5):1043--1065, 1996.
\bibitem[Welch(1967)]{welch1967} P.~D. Welch. The use of fast Fourier transform for the estimation of power spectra. \emph{IEEE Trans. Audio Electroacoustics}, 15(2):70--73, 1967.
\bibitem[Bazett(1920)]{bazett1920interval} H.~C. Bazett. An analysis of the time-relations of electrocardiograms. \emph{Heart}, 7:353--370, 1920.
\bibitem[Fridericia(1920)]{fridericia1920systole} L.~S. Fridericia. The duration of systole in an electrocardiogram in normal humans. \emph{Acta Medica Scandinavica}, 53:469--486, 1920.
\bibitem[Brier(1950)]{brier1950} G.~W. Brier. Verification of forecasts expressed in terms of probability. \emph{Monthly Weather Review}, 78(1):1--3, 1950.
\bibitem[DeLong et~al.(1988)]{delong1988} E.~R. DeLong, D.~M. DeLong, and D.~L. Clarke-Pearson. Comparing the areas under two or more correlated ROC curves: a nonparametric approach. \emph{Biometrics}, 44(3):837--845, 1988.
\bibitem[Vickers and Elkin(2006)]{vickers2006dca} A.~J. Vickers and E.~B. Elkin. Decision curve analysis: a novel method for evaluating prediction models. \emph{Medical Decision Making}, 26(6):565--574, 2006.
\bibitem[Lundberg and Lee(2017)]{lundberg2017shap} S.~M. Lundberg and S.-I. Lee. A unified approach to interpreting model predictions. In \emph{NeurIPS}, 2017.
\bibitem[Selvaraju et~al.(2017)]{selvaraju2017gradcam} R.~R. Selvaraju et~al. Grad-CAM: Visual explanations from deep networks via gradient-based localization. In \emph{ICCV}, 2017.
\bibitem[Platt(1999)]{platt1999} J.~Platt. Probabilistic outputs for SVMs and comparisons to regularized likelihood methods. In \emph{NIPS Workshop}, 1999.
\bibitem[Zadrozny and Elkan(2002)]{zadrozny2002calibration} B.~Zadrozny and C.~Elkan. Transforming classifier scores into accurate multiclass probability estimates. In \emph{KDD}, 2002.
\bibitem[Youden(1950)]{youden1950index} W.~J. Youden. Index for rating diagnostic tests. \emph{Cancer}, 3(1):32--35, 1950.
\bibitem[Steyerberg(2019)]{steyerberg2019clinical} E.~W. Steyerberg. \emph{Clinical Prediction Models}. Springer, 2019.
\bibitem[Riley et~al.(2020)]{riley2020sample} R.~D. Riley et~al. Calculating the sample size required for developing a clinical prediction model. \emph{BMJ}, 368:m441, 2020.
\bibitem[Collins et~al.(2015)]{collins2015tripod} G.~S. Collins et~al. Transparent reporting of a multivariable prediction model for individual prognosis or diagnosis (TRIPOD). \emph{Annals of Internal Medicine}, 162(1):55--63, 2015.
\end{thebibliography}

\end{document}